{{دیوان
| عنوان = پیامِ صبح
| ماخذ = ویکی ماخذ
| ماخذ_لنک = https://ur.wikisource.org/wiki/%D9%BE%DB%8C%D8%A7%D9%85%D9%90_%D8%B5%D8%A8%D8%AD
| مقابلہ = اقبال سائبر لائبریری (اقبال اکادمی پاکستان)
| مقابلہ_لنک = https://www.iqbalcyberlibrary.net
}}
=== (ماخوذ از لانگ فیلو) ===

<poem>
اُجالا جب ہُوا رخصت جبینِ شب کی افشاں کا
نسیمِ زندگی پیغام لائی صبحِ خنداں کا

جگایا بُلبلِ رنگیں نوا کو آشیانے میں
کنارے کھیت کے شانہ ہِلایا اُس نے دہقاں کا

طلسمِ ظلمتِ شب سُورۂ والنُّور سے توڑا
اندھیرے میں اُڑایا تاجِ زر شمعِ شبستاں کا

پڑھا خوابیدگانِ دَیر پر افسونِ بیداری
برہمن کو دیا پیغام خورشیدِ دُرَخشاں کا

ہُوئی بامِ حرم پر آ کے یوں گویا مؤذّن سے
نہیں کھٹکا ترے دل میں نمودِ مہرِ تاباں کا؟

پُکاری اس طرح دیوارِ گلشن پر کھڑے ہو کر
چٹک او غنچہ گُل! تُو مؤذّن ہے گُلستاں کا

دیا یہ حکم صحرا میں چلو اے قافلے والو!
چمکنے کو ہے جُگنو بن کے ہر ذرّہ بیاباں کا

سوئے گورِ غریباں جب گئی زندوں کی بستی سے
تو یوں بولی نظارا دیکھ کر شہرِ خموشاں کا

ابھی آرام سے لیٹے رہو، مَیں پھر بھی آؤں گی
سُلادوں گی جہاں کو، خواب سے تم کو جگاؤں گی
</poem>
